% rosetta/front.tex — 扉页与使用说明
\begin{titlepage}
\centering
\vspace*{2.2cm}
{\Huge\bfseries 数学 $\leftrightarrow$ 物理\par}
\vspace{0.4cm}
{\Huge\bfseries 罗塞塔词典\par}
\vspace{1.0cm}
{\Large 读 Kapustin--Witten《电磁对偶与几何 Langlands 纲领》用的双向对照手册\par}
\vspace{1.6cm}
{\large 从定义、例子、对应机制到直觉与常见误解\par}
\vspace{0.5cm}
{\large 112 条词条（103 条主干 + 4 条深水区 + 5 条读后修补） \quad 8 条背景术语 \quad 7 条轨道 \quad 3 张速查表 \quad 1 张逻辑流程图\par}
\vspace{2.4cm}
{\large 底本：A.~Kapustin, E.~Witten,\\
\emph{Electric--Magnetic Duality and the Geometric Langlands Program},\\
arXiv:hep-th/0604151v3（231 页）\par}
\vfill
{\large 2026 年 11 月\par}
\end{titlepage}

\section*{怎么用这本词典}
\addcontentsline{toc}{section}{怎么用这本词典}

KW06 的难处不在公式长度，而在**两套语言同时在场**：同一段论证里，物理侧说「$S$-对偶把 A-膜换成 B-膜」，数学侧说「Hitchin 模空间的复结构 $J$ 与辛结构 $K$ 互换」。读者若只熟悉一侧，会在另一侧的术语上反复卡住。

这本词典按**轨道**组织，每条词条固定七段：数学侧的定义与例子、物理侧的定义与例子、对应规则、直觉、易错。三种读法：

\begin{enumerate}[leftmargin=2em,itemsep=3pt]
  \item \textbf{顺读}：从 T1 到 T6，沿着译本的章节顺序建立整幅地图。
  \item \textbf{查读}：读原文遇到陌生词，直接在目录里找该词条（112 条按主题分组，词条标题为「中文（English）」；读后修补见 T8）。
  \item \textbf{对照读}：先看附录 A 的总表（一行一对），再挑最不熟悉的十几条精读。
\end{enumerate}

\textbf{约定。}「数学侧/物理侧」用蓝/橙标签区分；「对应」说明谁映到谁、自由参数怎么变；「直觉」允许类比但标注其边界；「易错」只收「按字面理解会推出错误结论」的坑，不收译名之争。词条末尾给出原文锚点（节号与印刷页码），便于翻回 KW06 核对。

\textbf{与译本的关系。}本词典是 KW06 中译注本的配套导航，术语与译本 `GLOSSARY` 完全一致（如 \emph{canonical coisotropic brane} = 典范余各向同性膜，\emph{eigenbrane} = 特征膜）。译文本身在另一册；本册只做对照，不重复译文。

\medskip
\noindent{\small\color{gcol}版权：原著版权归原作者；本词典为学习与研究用途的原创整理，文中数学陈述以原文为准。}
\clearpage
